\chapter{Who Defines the Versions}\label{ch:versions}

Once a work exists as a family of realizations, control over the family becomes a form of power. Someone decides which realizations exist, which are shown in a given place, who may see which, and at what price. In ordinary cinema these decisions are dispersed and mostly visible: a studio releases a cut, a censor classifies it, a theater programs it. In selective cinema they can be made inside a single screening, and they can be made invisibly. This chapter asks who should make them.

\section{Authorial and distributive variation}

It helps to distinguish two sources of variation. \emph{Authorial} variation is created by the makers of the work as part of its design: a director who makes comic and dramatic realizations, writers who compose a family of endings. \emph{Distributive} variation is introduced after the work is made, by those who handle it: a distributor who replaces a sequence for a particular market, a theater that disables a realization, a sponsor who pays for a stream with altered product placement, a regulator who requires a cut.

Authorial variation is part of the work. Distributive variation changes it, and it may do so in ways the makers never intended and the audience never learns about. The apparatus cannot tell the two apart; a stream is a stream. The difference lies in who introduced the variation and whether its introduction is declared.

A selective work should therefore carry its provenance. Which realizations were made by its authors, which were added or altered afterward, by whom, and why. A realization the authors did not make is not thereby illegitimate, since local realization in \Cref{ch:local} is often exactly this, but it should not be presented as the authors' work.

\section{Assignment by category}

The earliest illustration of selective cinema was an adult and a children's version on one screen, and it suggests assigning realizations by category: by age, and by extension by language, ticket class, region, disability, subscription tier, or inferred sensitivity. Some of these are reasonable. A viewer who needs a language will want the realization in it. Others amount to sorting people without telling them.

The book's position, developed across \Cref{ch:depth,ch:choice,ch:instruction}, is that silent assignment should not be the default. Viewers should be able to see what dimensions a family varies along, choose among them, and know when a choice has been made for them and on what basis. ``The version suitable for you'' is not an adequate description when the grounds of suitability are withheld. A viewer assigned a realization because of a category they belong to, or one a system inferred, is owed the knowledge that this happened.

\section{Price}

A family's realizations could be priced differently. The expert realization might cost more, a preferred ending might be sold as an upgrade, a local performance offered at a premium. Some of this may reflect real costs. Much of it would be artificial scarcity. Once several realizations are already being projected in a screening, every viewer's glasses could admit any of them; disabling one for a cheaper ticket is a decision about access, not a technical necessity.

Tiered pricing inside a single room also produces something ordinary cinema rarely does: visible class difference within one audience's experience of one work. Two people in adjacent seats, one of whom paid to see the ending and one of whom did not, are not sharing an occasion in the sense of \Cref{ch:synchronized-difference}. They are receiving a stratified product. A form whose promise is shared presence with different experiences should be wary of turning the difference into a price list.

\section{Honest description}

\Cref{ch:production} observed that a family's count of realizations is a poor measure of what it offers. Three dimensions of two options each make eight nominal realizations, whether they differ greatly or hardly at all. Advertising ``eight films in one'' when four of the eight differ only in subtitles misleads, and the temptation to inflate counts will be strong once the format is sold on them.

Honest description names the dimensions and says how far apart the realizations are along each. A family might be described as two genres of one story, sharing all events and differing in performance and score, with two endings that diverge in the final ten minutes, and captioned throughout. That is longer than a number and much more useful. It tells a viewer what they are choosing among.

\section{Labor and consent}

A variant family is made by more people doing more kinds of work than its single title suggests: writers of each column and of the spine, performers who play the same blocking in several registers, composers of a complementary score, editors balancing temporal debt, the continuity supervisor for the family proposed in \Cref{ch:production}, translators and local performers, the makers of access realizations. Credit should record who made each realization and who designed the architecture that holds them together. Otherwise the apparent unity of the work erases the plurality of its makers.

Performers raise a specific issue. An actor who agrees to play a scene in a thriller has not thereby agreed to have the performance appear in a satire, recut into a different context, or re-rendered to speak lines in another language. The flexibility that makes variant families possible, modular footage, re-editing, synthetic re-rendering, makes it easy to use a performance in realizations the performer never approved. Consent should be given per realization, and a performance should not appear in a realization, or a composite, that its performer has not agreed to.

\section{Control of the dial}

All of these questions reduce to one: who controls the dial? In the design this book has argued for, the viewer does. The glasses are timing devices set by hand; the realizations are declared; the variation is authorial or its provenance is stated; assignment, when it happens, is disclosed. Every departure from that design moves control from the viewer to someone else. Some departures are justified. None should be made silently. The next two chapters examine the two most consequential departures: glasses that only licensed viewers can use, and variation used to tell different audiences different things.
